\chapimg{fig/sh07.jpg}
\IfFileExists{extra/lead07.tex}{\chapterlead{\subhd{After this chapter you can} name the layers of the SWAFL application and state what each one computes, from which inputs and with which equations; reproduce one number of every layer by hand; and chain the layers into a complete design for a farm, from the soil and crop to the pump.}
}{}
\chapter{The SWAFL Layers: A Guided Tour}\label{chap:07}
\markboth{The SWAFL Layers: A Guided Tour}{The SWAFL Layers: A Guided Tour}
\IfFileExists{extra/th07.tex}{\section{What a decision-support system is}
A \kw{decision-support system} (DSS) is a program that turns data and a model into a recommendation that a manager can act on. It has three parts: a \emph{database} (weather, soil, crop, the state of the field), a \emph{model base} (the equations that compute the water balance, the design or the schedule) and a \emph{user interface} (the screens, the maps, the messages). The difference between a model and a DSS is the last step: the model gives numbers, the DSS gives advice (``irrigate in three days, 28 mm'') and says how sure it is.

\section{The architecture of SWAFL}
SWAFL is a set of \emph{layers} on one project. A layer is a page that answers one question (when to irrigate, how to design a lateral, where the water goes after a mulch). The layers use two \emph{engines}: Engine~1 (the annual meta-model, scaled to the results of the ensemble) and Engine~2 (the daily soil-water balance). The whole system is written so that the same project file can be opened by the online workbench and by the desktop application.
\begin{exlist}
\item \kw{The project.} It holds the field, the soil, the weather, the crop and the management, and everything that the layers compute; a change in one layer is seen in the others.
\item \kw{The engines.} They give the numbers; they are verified against Daisy (Chapter~\ref{chap:08}).
\item \kw{The layers.} Each combines the engines with the equations of one topic of this book.
\item \kw{The Daisy workbench.} It prepares, runs and reads Daisy simulations from the same project.
\end{exlist}

\section{Verification of a software}
A program that gives wrong answers with great confidence is dangerous. The software of this book is therefore tested at three levels: \emph{unit tests} (each equation against a hand calculation or a published example), \emph{integration tests} (the layers against each other: the sum of the parts equals the total) and \emph{system tests} (the answers against Daisy on held-out runs). A bug found at one level is corrected, and the test is added to the set, so that it cannot return.

\section{How to read a figure or a message of a layer}
Every figure of the layers is drawn from numbers, and the numbers come from equations of this book. To read a figure: (1) read the title and the axes, with their units; (2) find the equation that is behind it (the tour below names the chapter); (3) see whether the budget closes, since a figure whose terms do not add up is not to be trusted; (4) look for the warning messages, which say when a value is outside the range for which the equation was verified.
}{}
\IfFileExists{extra/bg07.tex}{\section{Introduction: the idea of a layer}
The SWAFL application (Soil Water, Atmosphere, Feedback and Layout) is organised in \emph{layers}. A layer is a self-contained view of one question about a field: when to irrigate, what a mulch does, how the structure of the soil changes drainage, how a plant grows, how the water balance of the field closes. All layers read the same inputs (the soil, the crop, the weather, the management) and write to a shared summary, so that a change in one place reaches all of them. The layers belong to two groups: those that \emph{read} the archived results of the ten-thousand Daisy simulations (control, mulch, structure, loops) and those that \emph{compute} with an engine (the field balance, the seasonal data, the crops, the designs, the phenotyping). The application also contains the \emph{Daisy workbench}, in which a user plans and reads a Daisy experiment. This chapter tours the layers in the order of the project tree; each section gives the background of the question, the main relation the layer uses, and one number reproduced by hand. The following chapters take each family to depth.

\section{Irrigation design: surface, drip and sprinkler}
The design layer takes a soil, a crop and a field and returns a system that can irrigate it: for surface irrigation the inflow, the advance and the cut-off time; for drip the emitter, the lateral and the wetted bulb; for sprinklers the nozzle, spacing and pressure; for all of them the main pipe and the pump. Its first relation is the \kw{system capacity}: to apply a net depth $d$ over an area $A$ in $T$ days of $h$ working hours each, with application efficiency $E_a$,

\begin{equation}Q=\frac{10\,A\,d}{E_a\,T\,h}\quad(\mathrm{m^3\,h^{-1}};\ A\ \text{in ha},\ d\ \text{in mm}).\label{eq:cap7}\end{equation}
\par\noindent where $Q$ is the system capacity (m$^3$\,h$^{-1}$), $A$ the area irrigated (ha), $d$ the gross depth per irrigation (mm), $E_a$ the application efficiency, $T$ the irrigation interval (d) and $h$ the hours of operation per day.



\exbox{System capacity}{A 40-ha maize field must receive a net 45 mm in 10 days, working 20 h per day, with sprinklers of $E_a=0.75$. Find the system capacity in m$^3$\,h$^{-1}$ and in L\,s$^{-1}$.}\sol \givenline{$E_a=0.75$}\ans \[Q=10\times40\times45/(0.75\times10\times20)=120\] m$^3$\,h$^{-1}$ $=33.3$ L\,s$^{-1}$ \ans That is the flow to be carried by the mainline and delivered by the pump (Chapter~7); at 60 L\,s$^{-1}$ per block it needs 0.6 blocks working at once. \finans

\section{Irrigation control: sensor and trigger}
The control layer plays a season of soil water and shows when a sensor at a chosen depth passes the trigger. The soil column is a daily model; the sensor reads the suction from the retention curve; the trigger is a suction or a depletion. The numbers that come out are the days beyond the trigger, the irrigation applied and the yield, taken for each soil and trigger from the archived experiments (Chapter~4). A simple version of the same logic can be run by hand with a depletion model.

\exbox{Count of irrigations}{A crop uses 5 mm\,d$^{-1}$ for 100 days from a root zone with RAW = 30 mm refilled at each irrigation. How many irrigations are triggered, and what is the net water applied?}\sol The interval is $30/5=6$ d; irrigations on days 6, 12, \dots, 96: $\lfloor100/6\rfloor=16$ irrigations \ans Net water $=16\times30=480$ mm \ans The remaining 4 days draw the store down by 20 mm, which is the end-of-season depletion. \finans

\section{Mulch and evaporation}
A mulch lowers the vapour flux from the soil by the \emph{strength} $s$ (the fraction by which the flux is reduced) and holds a depth of rain, the \emph{capacity} $c$, that evaporates from the mulch itself. The layer reads the Daisy experiments over strength, rain and capacity and shows the evaporation, drainage, irrigation and yield of the mulched run against the bare one, and the fate of the water saved (Chapter~10). To first order, the evaporation from a mulched soil is $E_m\approx(1-s)E_s+E_{int}$, with $E_{int}=\sum\min(P_j,c)$ over the rain events.

\exbox{Evaporation from a mulched soil}{Bare-soil evaporation is 180 mm per season. A mulch with strength 0.6 and capacity 2 mm meets 35 rain events of more than 2 mm. Estimate the evaporation of the mulched soil and the apparent saving.}\sol \[E_m=(1-0.6)\times180+35\times2=72+70=142\ \text{mm}\]
apparent saving $180-142=38$ mm \ans The interception of the mulch gives back 65 per cent of the saving of 108 mm, which is why a mulch of large capacity saves less than its strength suggests. \finans

\section{Soil structure and tillage}
Tillage loosens the topsoil (lowers the bulk density, raises the porosity and the conductivity) and the soil then consolidates with time and rain. The layer shows the effect of tillage (conventional, reduced, none) and of the consolidation rate on drainage, evaporation and yield (Chapter~10). A simple model: the bulk density relaxes exponentially, $\rho_b(t)=\rho_\infty-(\rho_\infty-\rho_0)e^{-kt}$, and the saturated conductivity scales with the Kozeny--Carman factor $K_s\propto\phi^3/(1-\phi)^2$.

\exbox{Consolidation after tillage}{After ploughing, $\rho_b=1.15$ Mg\,m$^{-3}$ and it relaxes towards $1.45$ with $k=0.02$ d$^{-1}$. Find $\rho_b$ at 0, 30, 60, 120 and 240 days and the relative saturated conductivity $K_s(t)/K_s(0)$ ($\rho_s=2.65$).}\sol \givenline{$\rho_b=1.15$ Mg\,m$^{-3}$ ; $k=0.02$ d$^{-1}$ ; $\rho_s=2.65$}\ans $t=0$: $\rho_b=1.150$, $K_s/K_{s0}=1.000$; $t=30$: $\rho_b=1.285$, $K_s/K_{s0}=0.603$; $t=60$: $\rho_b=1.360$, $K_s/K_{s0}=0.455$; $t=120$: $\rho_b=1.423$, $K_s/K_{s0}=0.358$; $t=240$: $\rho_b=1.448$, $K_s/K_{s0}=0.325$ \ans The conductivity of the tilled layer falls to about a third within the first two months, which is the reason the benefit of tillage on drainage is short. \finans

\section{Phenotyping: the deterministic plant and the volumetric scan}
The two phenotyping layers grow a plant by a deterministic crop model and then scan it with a virtual camera rig to measure its traits (Chapter~11). The traits are height, crown size, volume, leaf area and biomass. A convenient geometric model of a plant is an ellipsoid of semi-axes $a$, $a$ (horizontal) and $b$ (vertical): volume $V=\frac43\pi a^2b$, and the leaf area index of the ground area $A_g=\pi a^2$ is $\mathrm{LAI}=\mathrm{LA}/A_g$ with LA the leaf area.

\exbox{Volume and LAI of a plant}{A plant of height $H=0.9$ m has a crown of horizontal radius 0.35 m and is modelled as an ellipsoid. Find its volume and, if its leaf area is 0.85 m$^2$, the LAI of the ground area it covers.}\sol \givenline{$H=0.9$ m has a crown}\ans Semi-axes $a=0.35$ m, $b=H/2=0.45$ m: \[V=\frac43\pi\times0.35^2\times0.45=0.231\ \text{m$^3$}\] \ans Ground area $A_g=\pi a^2=0.385$ m$^2$, so \[\mathrm{LAI}=0.85/0.385=2.21\quad\text{\finans}\]

\section{Field water balance}
The field-balance layer runs the daily bucket (Chapter~9): the depletion $D_r$ of the root zone is updated each day, $D_{r,i}=D_{r,i-1}-P_i-I_i+ET_{a,i}+D_{p,i}$, with deep percolation $D_{p}=\max(0,P+I-ET_a-D_{r,i-1})$ when the soil is refilled beyond field capacity, and $ET_a=K_sK_cET_0$.

\begin{table}[!htbp]\centering\small\caption{One week of the daily bucket: TAW = 120 mm, RAW = 60 mm, $K_c=1.1$, initial depletion 40 mm; all in mm.}
\begin{tabular}{rrrrrr}\toprule
\hd{Day} & \hd{$ET_0$} & \hd{Rain} & \hd{$K_s$} & \hd{$ET_a$} & \hd{$D_r$ end}\\\midrule
1 & 5 & 0 & 1.00 & 5.50 & 45.5\\
2 & 5.5 & 0 & 1.00 & 6.05 & 51.5\\
3 & 6 & 0 & 1.00 & 6.60 & 58.1\\
4 & 6 & 12 & 1.00 & 6.60 & 52.8\\
5 & 5.5 & 0 & 1.00 & 6.05 & 58.8\\
6 & 5 & 0 & 1.00 & 5.50 & 64.3\\
7 & 4.5 & 0 & 0.93 & 4.60 & 68.9\\\bottomrule\end{tabular}\end{table}

Day by day the depletion grows by the actual evapotranspiration; the rain of 12 mm on day 4 lowers it; and as $D_r$ approaches RAW the stress factor $K_s$ starts to fall below 1. Seven days after the start the depletion is 68.9 mm: the crop is beyond the readily available water.

\section{Seasonal data and projections}
The seasonal layers summarise a year of weather and the water balance as charts and tables, and project them forward by a rate of warming and of rainfall change per decade. A projection to the year $y$ after the record applies $\Delta T=r_T\,y/10$ and $P\to P(1+r_P\,y/10)$; the reference evapotranspiration rises by about 2 to 4 per cent per $^\circ$C of warming in humid climates. The aridity index $AI=P/ET_0$ classifies the climate: below 0.2 is arid, 0.2--0.5 semi-arid, 0.5--0.65 dry sub-humid, above 0.65 humid.

\exbox{Aridity index and projection}{A site has $P=560$ mm and $ET_0=1\,120$ mm. The projection is $+0.3$ $^\circ$C and $-2$ per cent rainfall per decade. Find the aridity index now and after 30 years if $ET_0$ rises 3 per cent per $^\circ$C.}\sol \givenline{$P=560$ mm ; $ET_0=1\,120$ mm}\ans Now: $AI=560/1120=0.50$ (dry sub-humid, at the limit with semi-arid). After 30 years: $\Delta T=0.9$ $^\circ$C; \[ET_0=1120\times(1+0.03\times0.9)=1150\ \text{mm}\]
\[P=560\times(1-0.02\times3)=526\ \text{mm}\]
$AI=0.458$ \ans The climate moves to the semi-arid class, and the net irrigation requirement rises by about 8 per cent for a crop of $ET_c\approx0.9ET_0$. \finans

\section{Crops and the crop water requirement}
The crop layer holds the library of 48 crops (Chapter~6) and computes, for a crop, a soil and a climate, the seasonal $ET_c$, the net and gross irrigation requirement, the water productivity and the water footprint. The water productivity is $WP=Y/ET_a$ (kg of yield per m$^3$ of water evapotranspired); typical values are 1.2--2.0 kg\,m$^{-3}$ for maize grain and 0.8--1.2 for wheat.

\exbox{Water productivity}{Maize yields 10.8 Mg\,ha$^{-1}$ with $ET_a=560$ mm; wheat yields 6.4 Mg\,ha$^{-1}$ with $ET_a=420$ mm. Compare their water productivities.}\sol \givenline{$ET_a=560$ mm; wheat yields ; $ET_a=420$ mm}\ans Maize: $WP=10\,800/5\,600=1.93$ kg\,m$^{-3}$; wheat: $6\,400/4\,200=1.52$ kg\,m$^{-3}$ \ans Maize gets more grain per cubic metre of water, as a C4 crop does. \finans

\section{Feedback loops}
The loops layer draws the links between climate, management, soil water, crop and nitrogen as arrows, and runs the Daisy results around them: more irrigation raises the soil water, which raises transpiration and yield but also drainage and nitrogen leaching; a mulch lowers evaporation, which raises the store and so drainage. A loop has a sign (positive if it reinforces, negative if it balances) equal to the product of the signs of its arrows. The nitrogen lost by leaching is $L_N=D\,c_N$ with $D$ the drainage and $c_N$ the concentration of nitrate in the drainage water.

\exbox{Nitrogen lost with drainage}{An extra 120 mm of irrigation raises drainage by 30 mm. The drainage water carries 18 mg\,L$^{-1}$ of nitrate-nitrogen. Find the extra nitrogen lost per hectare and the sign of the loop irrigation $\to$ drainage $\to$ nitrogen loss $\to$ yield.}\sol \[L_N=30\,\mathrm{mm}\times18\,\mathrm{mg\,L^{-1}}=30\,\mathrm{L\,m^{-2}}\times0.018\,\mathrm{g\,L^{-1}}=0.54\] g\,m$^{-2}=5.4$ kg\,ha$^{-1}$ \ans Signs: irrigation $\to$ drainage $(+)$; drainage $\to$ nitrogen loss $(+)$; nitrogen loss $\to$ yield $(-)$; the product is negative, so the loop is a \emph{balancing} loop: the extra water that the crop does not use removes the nitrogen that it needs. \finans

\section{The Daisy model in the application}
The last branch of the project tree is the Daisy workbench, with which a user chooses a Daisy experiment from a ready list or builds one, runs the combinations, reads the results in tables and charts and exports them (Chapters~3 and~9). The experiments of this book are read from the archive of 10\,190 runs; the workbench itself runs on any Daisy installation. The workbench teaches the \emph{logic of an experiment}: a base case, a set of factors with levels, the full factorial of their combinations, and a response. For $k$ factors with $L_i$ levels the number of runs is $\prod L_i$; main effects are the differences of the mean response between levels, and interactions are the departures from additivity.

\exbox{}{A design has three factors: soil (3 levels), mulch strength (3 levels) and rain scale (2 levels). Find the number of runs. In the table of mean yield (Mg\,ha$^{-1}$): mulch 0 = 6.1, mulch 0.5 = 6.5, mulch 1.0 = 6.8; find the main effect of the strongest mulch against none and its relative size.}\sol Runs $=3\times3\times2=18$ \ans Main effect $=6.8-6.1=0.7$ Mg\,ha$^{-1}$ \ans Relative $=0.7/6.1=11.5$ per cent \finans

}{}
\banner{A tour of the SWAFL layers}
\section{The workbench}
\fig{sh_full}{0.98}{The SWAFL workbench: the layer tree on the left, the simulation window, the properties and the dashboard. The date in the toolbar moves the whole workbench.}
There are 32 layers: 23 for decisions, in nine groups, and nine that work with Daisy. Version 1.1.0 adds seasons projected 10 or 20 years ahead, a pipe network with Hazen--Williams losses and pump head, a drip wetting pattern that follows the soil, and a Daisy scenario designer. Weather comes from a built-in record, your own file, or an online archive.

\section{Field balance and the soil--plant--air column}
\fig{sc_fieldk}{0.8}{Field water balance: precipitation, irrigation, evapotranspiration, drainage and storage (Equation~\ref{eq:bal}).}
\fig{sc_spack}{0.8}{Soil, plant and air: the water potential of each compartment, from the soil particles to the atmosphere.}
\exbox{}{Open the field balance for the sand and for the silty clay loam. Which has the larger drainage? Relate your answer to Table~\ref{tab:soils}.}\sol  Sand (S03) has the larger drainage. Table~\ref{tab:soils} gives the memory fate $\phi_D=0.846$ for the sand and $0.022$ for the silty clay loam: of each millimetre of stored surplus the sand drains $85$ per cent and the silty clay loam $2$ per cent. The ratio is $0.846/0.022=38$. \finans


\section{Mulch, evaporation and the fate of saved water}
\fig{sc_n2k}{0.8}{Fate of the water saved by mulch across rainfall. The drained share depends on the soil (Chapter~8).}
\fig{sc_n3k}{0.8}{The regime map: where vapour gives way to drainage.}
\exbox{}{Use the competition law of Chapter~8 to predict, for sand and for silty clay loam, which of the two soils shows more drainage in the fate-of-saved-water layer, and check on the screen.}\sol  Invert $\phi_D=\Ks/(\Ks+q^*)$ to get $\Ks/q^*=\phi_D/(1-\phi_D)$. Sand: $0.846/0.154=5.5$. Silty clay loam: $0.022/0.978=0.0225$. The conductivity-to-withdrawal ratio of the sand is $5.5/0.0225\approx 244$ times larger, so the sand drains far more of the saved water; on the screen the drainage bar of the sand dominates and that of the silty clay loam is nearly empty. \finans


\section{Triggers, structure and roots}
\fig{sc_v1k}{0.8}{Sensor suction every day from 2019 to 2039. The dashed line is the trigger. Irrigated (blue) against rainfed (orange); bars give the growing-season days beyond the trigger. In 2029 the irrigated field had no days beyond 471 cm and the rainfed field had 101 of 183.}
\fig{sc_n1k}{0.8}{The trigger law: the depletion of the root zone at the trigger is soil independent.}
\fig{sc_r1k}{0.8}{The root layer: where the water is taken.}
\exbox{}{In the sensor layer set the trigger of the sand to 300 cm. Predict, with Table~\ref{tab:reach}, what you will see, then check.}\sol  For the sand (S03), $\hc=199$ cm. The ratio is $300/199=1.51>1$, so the trigger lies beyond the reach and cannot fire. The archived trigger runs agree: mean irrigation is $122$, $38.5$ and $6.4$ mm yr$^{-1}$ at 50, 75 and 100 cm, and $0$ from 150 cm upward. On the screen you see an irrigation total of zero with the crop on rainfall alone. \finans


\section{Irrigation design: surface, drip, sprinkler}
\exbox{}{In the surface layer compare the advance time for the same inflow on a loam and a sand, and explain the difference from the intake.}\sol  Use the volume balance $Q_0t=\sigma_zLZ(t)$ with $\sigma_z=0.7$, $L=100$ m, $Q_0=0.001$ m$^2$\,s$^{-1}$ per unit width and the Kostiakov intake $Z=k\tau^a$ (cm, $\tau$ in min). The parameters below are \emph{illustrative}, not measured: sand $k=1.2$, $a=0.6$; loam $k=0.5$, $a=0.45$. Sand: $0.06t=70(0.012)t^{0.6}$ gives $t^{0.4}=14$, so $t=733$ min. Loam: $0.06t=70(0.005)t^{0.45}$ gives $t^{0.55}=5.83$, so $t=24.7$ min. The sand takes thirty times longer to advance (it never completes the run in a realistic furrow) because its intake is much larger at all times. \finans


\section{Drainage and crops}
\fig{sc_cropk}{0.8}{Crop coefficients, daily water need and the irrigation schedule. The library has 75 crops and follows the FAO-56 edition of 2025.}

\section{Seasons, balances and feedback loops}
\fig{sc_wb2k}{0.8}{Full water balance year by year. Seasons after 2038 are projections.}
\fig{sc_nbk}{0.8}{Nitrogen budget year by year.}
\fig{sc_v5k}{0.8}{Feedback loops between climate, management, soil water, crop and nitrogen. The red arrows join the parts that act on each other; the table lists what one step of each control does.}

\section{Phenotyping and the real field}
\fig{sc_dphk}{0.75}{Deterministic phenotyping: an exactly reproducible virtual experiment whose water balance closes (set-up form shown).}
\fig{sc_vphk}{0.75}{Volumetric phenotyping: a virtual three-dimensional scanner that measures volume, crown, leaves and height.}
\exbox{}{Explain why an exactly reproducible virtual experiment is useful for teaching, and name one limit.}\sol  A virtual experiment can be repeated exactly: the same weather, soil and settings give the same numbers, so a result can be checked by anyone and one factor changed at a time. One limit: the model is one-dimensional matrix flow, so macropores, lateral flow and real field variability are absent, and no field validation is implied.
\IfFileExists{extra/app07.tex}{\subsection{The output of the application, panel by panel}
The layer \emph{Soil, plant and air} is the picture of the whole chain that Chapter 8 describes in words: one continuous water column from the groundwater to the atmosphere, in which the suction rises from $-0.02$ MPa in the soil to $-56$ MPa in the air.

Each height of the column from the water table to the leaves has its own suction and content; the arrows show the direction of the flow.

\appfig{spac_0}{0.96}{Soil to sky: one continuous water column, with the suction, the water content and the pores holding water at each height.}

The figure compares the pull of the atmosphere on the water with the hold of the soil on it and shows that the water moves along the gradient of the total head.

\appfig{spac_1}{0.96}{The air pulls with $-56$ MPa; the soil holds on with $-0.02$ MPa.}

The three profiles are the simulated state of the column on the day of the replay.

\appfig{spac_2}{0.96}{Simulated suction, water content and pores holding water along the column.}

}{}
\IfFileExists{extra/proj07.tex}{\banner{Design project 7: a complete design for a 40-ha maize farm, layer by layer}
\textbf{Brief.} A farm of 40 ha on loam (S06) grows maize. The weather is that of a semi-arid summer, with a peak $ET_0$ of 7 mm\,d$^{-1}$. The farmer asks for a sensor-controlled sprinkler system with a trigger, a pipe and a pump. This project chains the layers of the application: the crop and weather layer gives the demand, the soil and control layers the dose and the interval, the design layer the capacity, and the pipe and the pump follow. Every step of the table is computed by hand.

\begin{table}[!htbp]\centering\small\caption{The chain of the design: each line uses the result above it.}
\begin{tabular}{lll}\toprule
\hd{Layer} & \hd{Quantity} & \hd{Result}\\\midrule
Crop & $K_{c,mid}\,ET_{0,peak}$ & $1.20\times7.0=8.4$ mm/d\\
Soil & TAW $=1.7\times90$ & 153 mm\\
Control & RAW $=0.5\,$TAW & 76 mm\\
Control & Interval $=RAW/ET_c$ & 9.1 d $\to$ 9 d\\
Design & Gross depth $RAW/0.75$ & 102 mm\\
Design & Capacity $10Ad/(E_aTh)$ & 227 m$^3$/h = 63.0 L/s\\
Pipe & Velocity in $D=200$ mm & 2.00 m/s\\
Pipe & $h_f$ (800 m, $C=140$) & 13.7 m\\
Pump & Head $35+h_f$ & 48.7 m\\
Pump & Power $\rho gQH/0.65$ & 46.3 kW\\\bottomrule\end{tabular}\end{table}

\textbf{Reading.} The soil can store 76 mm of readily available water, which a crop at the peak demand of 8.4 mm\,d$^{-1}$ uses in 9.1 days; the cycle is therefore 9 days. A gross depth of 102 mm every 9 days on 40 ha, applied in 20 h per day, calls for 63.0 L\,s$^{-1}$. In a 200 mm mainline this flow runs at 2.00 m\,s$^{-1}$ (above the limit of 1.5) and loses 13.7 m of head, so the pump has to give 48.7 m and draws 46.3 kW.

\textbf{Sensitivity.} If the system efficiency falls from 0.75 to 0.65 (a worn sprinkler package), the capacity rises by $0.75/0.65-1=15$ per cent and the friction loss by about $(1.154)^{1.852}-1=30$ per cent, so the pump power grows by about 33 per cent: an efficiency lost in the field is paid for in the pump house.

}{}
\IfFileExists{extra/ten07.tex}{\section{Worked solutions on the ten soils}
The same 40-ha farm is sized on each of the ten soils: the soil fixes the interval and, with it, the capacity.

\exbox{Ten-soil example 7a: Coarse loam (S01)}{Coarse loam: size the irrigation interval and system capacity of a 40-ha maize farm at peak $ET_c=6.5$ mm\,d$^{-1}$, root zone 0.6 m, $p=0.5$, $E_a=0.75$, working 20 h per day.}\sol $TAW=90$ mm, $RAW=45$ mm \ans Interval $T=\lfloor RAW/ET_c\rfloor=\lfloor6.9\rfloor=6$ d \ans Net depth $T\,ET_c=39.0$ mm; gross $=39.0/0.75=52.0$ mm \ans Capacity $Q=10Ad/(E_aTh)=10\times40\times39.0/(0.75\times6\times20)=173$ m$^3$\,h$^{-1}$ ($48$ L\,s$^{-1}$) \ans

\exbox{Ten-soil example 7a: Coarse sand (S02)}{Coarse sand: size the irrigation interval and system capacity of a 40-ha maize farm at peak $ET_c=6.5$ mm\,d$^{-1}$, root zone 0.6 m, $p=0.5$, $E_a=0.75$, working 20 h per day.}\sol $TAW=30$ mm, $RAW=15$ mm \ans Interval $T=\lfloor RAW/ET_c\rfloor=\lfloor2.3\rfloor=2$ d \ans Net depth $T\,ET_c=13.0$ mm; gross $=13.0/0.75=17.3$ mm \ans Capacity $Q=10Ad/(E_aTh)=10\times40\times13.0/(0.75\times2\times20)=173$ m$^3$\,h$^{-1}$ ($48$ L\,s$^{-1}$) \ans

\exbox{Ten-soil example 7a: Sand (S03)}{Sand: size the irrigation interval and system capacity of a 40-ha maize farm at peak $ET_c=6.5$ mm\,d$^{-1}$, root zone 0.6 m, $p=0.5$, $E_a=0.75$, working 20 h per day.}\sol $TAW=36$ mm, $RAW=18$ mm \ans Interval $T=\lfloor RAW/ET_c\rfloor=\lfloor2.8\rfloor=2$ d \ans Net depth $T\,ET_c=13.0$ mm; gross $=13.0/0.75=17.3$ mm \ans Capacity $Q=10Ad/(E_aTh)=10\times40\times13.0/(0.75\times2\times20)=173$ m$^3$\,h$^{-1}$ ($48$ L\,s$^{-1}$) \ans

\exbox{Ten-soil example 7a: Loamy sand (S04)}{Loamy sand: size the irrigation interval and system capacity of a 40-ha maize farm at peak $ET_c=6.5$ mm\,d$^{-1}$, root zone 0.6 m, $p=0.5$, $E_a=0.75$, working 20 h per day.}\sol $TAW=54$ mm, $RAW=27$ mm \ans Interval $T=\lfloor RAW/ET_c\rfloor=\lfloor4.2\rfloor=4$ d \ans Net depth $T\,ET_c=26.0$ mm; gross $=26.0/0.75=34.7$ mm \ans Capacity $Q=10Ad/(E_aTh)=10\times40\times26.0/(0.75\times4\times20)=173$ m$^3$\,h$^{-1}$ ($48$ L\,s$^{-1}$) \ans

\exbox{Ten-soil example 7a: Sandy loam (S05)}{Sandy loam: size the irrigation interval and system capacity of a 40-ha maize farm at peak $ET_c=6.5$ mm\,d$^{-1}$, root zone 0.6 m, $p=0.5$, $E_a=0.75$, working 20 h per day.}\sol $TAW=72$ mm, $RAW=36$ mm \ans Interval $T=\lfloor RAW/ET_c\rfloor=\lfloor5.5\rfloor=5$ d \ans Net depth $T\,ET_c=32.5$ mm; gross $=32.5/0.75=43.3$ mm \ans Capacity $Q=10Ad/(E_aTh)=10\times40\times32.5/(0.75\times5\times20)=173$ m$^3$\,h$^{-1}$ ($48$ L\,s$^{-1}$) \ans

\exbox{Ten-soil example 7a: Loam (S06)}{Loam: size the irrigation interval and system capacity of a 40-ha maize farm at peak $ET_c=6.5$ mm\,d$^{-1}$, root zone 0.6 m, $p=0.5$, $E_a=0.75$, working 20 h per day.}\sol $TAW=102$ mm, $RAW=51$ mm \ans Interval $T=\lfloor RAW/ET_c\rfloor=\lfloor7.8\rfloor=7$ d \ans Net depth $T\,ET_c=45.5$ mm; gross $=45.5/0.75=60.7$ mm \ans Capacity $Q=10Ad/(E_aTh)=10\times40\times45.5/(0.75\times7\times20)=173$ m$^3$\,h$^{-1}$ ($48$ L\,s$^{-1}$) \ans

\exbox{Ten-soil example 7a: Silt loam (S07)}{Silt loam: size the irrigation interval and system capacity of a 40-ha maize farm at peak $ET_c=6.5$ mm\,d$^{-1}$, root zone 0.6 m, $p=0.5$, $E_a=0.75$, working 20 h per day.}\sol $TAW=120$ mm, $RAW=60$ mm \ans Interval $T=\lfloor RAW/ET_c\rfloor=\lfloor9.2\rfloor=9$ d \ans Net depth $T\,ET_c=58.5$ mm; gross $=58.5/0.75=78.0$ mm \ans Capacity $Q=10Ad/(E_aTh)=10\times40\times58.5/(0.75\times9\times20)=173$ m$^3$\,h$^{-1}$ ($48$ L\,s$^{-1}$) \ans

\exbox{Ten-soil example 7a: Sandy clay loam (S08)}{Sandy clay loam: size the irrigation interval and system capacity of a 40-ha maize farm at peak $ET_c=6.5$ mm\,d$^{-1}$, root zone 0.6 m, $p=0.5$, $E_a=0.75$, working 20 h per day.}\sol $TAW=84$ mm, $RAW=42$ mm \ans Interval $T=\lfloor RAW/ET_c\rfloor=\lfloor6.5\rfloor=6$ d \ans Net depth $T\,ET_c=39.0$ mm; gross $=39.0/0.75=52.0$ mm \ans Capacity $Q=10Ad/(E_aTh)=10\times40\times39.0/(0.75\times6\times20)=173$ m$^3$\,h$^{-1}$ ($48$ L\,s$^{-1}$) \ans

\exbox{Ten-soil example 7a: Clay loam (S09)}{Clay loam: size the irrigation interval and system capacity of a 40-ha maize farm at peak $ET_c=6.5$ mm\,d$^{-1}$, root zone 0.6 m, $p=0.5$, $E_a=0.75$, working 20 h per day.}\sol $TAW=96$ mm, $RAW=48$ mm \ans Interval $T=\lfloor RAW/ET_c\rfloor=\lfloor7.4\rfloor=7$ d \ans Net depth $T\,ET_c=45.5$ mm; gross $=45.5/0.75=60.7$ mm \ans Capacity $Q=10Ad/(E_aTh)=10\times40\times45.5/(0.75\times7\times20)=173$ m$^3$\,h$^{-1}$ ($48$ L\,s$^{-1}$) \ans

\exbox{Ten-soil example 7a: Silty clay loam (S10)}{Silty clay loam: size the irrigation interval and system capacity of a 40-ha maize farm at peak $ET_c=6.5$ mm\,d$^{-1}$, root zone 0.6 m, $p=0.5$, $E_a=0.75$, working 20 h per day.}\sol $TAW=102$ mm, $RAW=51$ mm \ans Interval $T=\lfloor RAW/ET_c\rfloor=\lfloor7.8\rfloor=7$ d \ans Net depth $T\,ET_c=45.5$ mm; gross $=45.5/0.75=60.7$ mm \ans Capacity $Q=10Ad/(E_aTh)=10\times40\times45.5/(0.75\times7\times20)=173$ m$^3$\,h$^{-1}$ ($48$ L\,s$^{-1}$) \ans

}{}
\IfFileExists{extra/tenb07.tex}{\subsection{Worked solutions on ten crops}
The ten crops of the library differ in their demand, their roots and their sensitivity. The same type of question is answered for each of them at $ET_0=5.5$ mm\,d$^{-1}$; the data are typical values of the crops (Doorenbos--Kassam and FAO 56).

\exbox{Crop example 7a: Maize}{Maize on 10 ha of the loam S06: design the interval, the gross depth and the flow to irrigate in 20 h\,d$^{-1}$ at $E_a=0.80$.}\sol \givenline{$E_a=0.80$}\ans $RAW=94$ mm; \[T=\lfloor94/6.60\rfloor=14\ \text{d}\] \ans Net depth $92.4$ mm; gross $115.5$ mm \ans \[Q=10Ad/(E_aTh)=10\times10\times92.4/(0.8\times14\times20)=41.2\] m$^3$\,h$^{-1}$ ($11.5$ L\,s$^{-1}$) \finans

\exbox{Crop example 7a: Wheat}{Wheat on 10 ha of the loam S06: design the interval, the gross depth and the flow to irrigate in 20 h\,d$^{-1}$ at $E_a=0.80$.}\sol \givenline{$E_a=0.80$}\ans $RAW=112$ mm; \[T=\lfloor112/6.32\rfloor=17\ \text{d}\] \ans Net depth $107.5$ mm; gross $134.4$ mm \ans \[Q=10Ad/(E_aTh)=10\times10\times107.5/(0.8\times17\times20)=39.5\] m$^3$\,h$^{-1}$ ($11.0$ L\,s$^{-1}$) \finans

\exbox{Crop example 7a: Cotton}{Cotton on 10 ha of the loam S06: design the interval, the gross depth and the flow to irrigate in 20 h\,d$^{-1}$ at $E_a=0.80$.}\sol \givenline{$E_a=0.80$}\ans $RAW=122$ mm; \[T=\lfloor122/6.32\rfloor=19\ \text{d}\] \ans Net depth $120.2$ mm; gross $150.2$ mm \ans \[Q=10Ad/(E_aTh)=10\times10\times120.2/(0.8\times19\times20)=39.5\] m$^3$\,h$^{-1}$ ($11.0$ L\,s$^{-1}$) \finans

\exbox{Crop example 7a: Tomato}{Tomato on 10 ha of the loam S06: design the interval, the gross depth and the flow to irrigate in 20 h\,d$^{-1}$ at $E_a=0.80$.}\sol \givenline{$E_a=0.80$}\ans $RAW=48$ mm; $T=\lfloor48/6.32\rfloor=7$ d \ans Net depth $44.3$ mm; gross $55.3$ mm \ans \[Q=10Ad/(E_aTh)=10\times10\times44.3/(0.8\times7\times20)=39.5\] m$^3$\,h$^{-1}$ ($11.0$ L\,s$^{-1}$) \finans

\exbox{Crop example 7a: Potato}{Potato on 10 ha of the loam S06: design the interval, the gross depth and the flow to irrigate in 20 h\,d$^{-1}$ at $E_a=0.80$.}\sol \givenline{$E_a=0.80$}\ans $RAW=30$ mm; $T=\lfloor30/6.32\rfloor=4$ d \ans Net depth $25.3$ mm; gross $31.6$ mm \ans \[Q=10Ad/(E_aTh)=10\times10\times25.3/(0.8\times4\times20)=39.5\] m$^3$\,h$^{-1}$ ($11.0$ L\,s$^{-1}$) \finans

\exbox{Crop example 7a: Soybean}{Soybean on 10 ha of the loam S06: design the interval, the gross depth and the flow to irrigate in 20 h\,d$^{-1}$ at $E_a=0.80$.}\sol \givenline{$E_a=0.80$}\ans $RAW=51$ mm; $T=\lfloor51/6.32\rfloor=8$ d \ans Net depth $50.6$ mm; gross $63.2$ mm \ans \[Q=10Ad/(E_aTh)=10\times10\times50.6/(0.8\times8\times20)=39.5\] m$^3$\,h$^{-1}$ ($11.0$ L\,s$^{-1}$) \finans

\exbox{Crop example 7a: Sugar beet}{Sugar beet on 10 ha of the loam S06: design the interval, the gross depth and the flow to irrigate in 20 h\,d$^{-1}$ at $E_a=0.80$.}\sol \givenline{$E_a=0.80$}\ans $RAW=84$ mm; \[T=\lfloor84/6.60\rfloor=12\ \text{d}\] \ans Net depth $79.2$ mm; gross $99.0$ mm \ans \[Q=10Ad/(E_aTh)=10\times10\times79.2/(0.8\times12\times20)=41.2\] m$^3$\,h$^{-1}$ ($11.5$ L\,s$^{-1}$) \finans

\exbox{Crop example 7a: Sunflower}{Sunflower on 10 ha of the loam S06: design the interval, the gross depth and the flow to irrigate in 20 h\,d$^{-1}$ at $E_a=0.80$.}\sol \givenline{$E_a=0.80$}\ans $RAW=61$ mm; \[T=\lfloor61/6.05\rfloor=10\ \text{d}\] \ans Net depth $60.5$ mm; gross $75.6$ mm \ans \[Q=10Ad/(E_aTh)=10\times10\times60.5/(0.8\times10\times20)=37.8\] m$^3$\,h$^{-1}$ ($10.5$ L\,s$^{-1}$) \finans

\exbox{Crop example 7a: Onion}{Onion on 10 ha of the loam S06: design the interval, the gross depth and the flow to irrigate in 20 h\,d$^{-1}$ at $E_a=0.80$.}\sol \givenline{$E_a=0.80$}\ans $RAW=18$ mm; $T=\lfloor18/5.78\rfloor=3$ d \ans Net depth $17.3$ mm; gross $21.7$ mm \ans \[Q=10Ad/(E_aTh)=10\times10\times17.3/(0.8\times3\times20)=36.1\] m$^3$\,h$^{-1}$ ($10.0$ L\,s$^{-1}$) \finans

\exbox{Crop example 7a: Alfalfa}{Alfalfa on 10 ha of the loam S06: design the interval, the gross depth and the flow to irrigate in 20 h\,d$^{-1}$ at $E_a=0.80$.}\sol \givenline{$E_a=0.80$}\ans $RAW=94$ mm; \[T=\lfloor94/6.32\rfloor=14\ \text{d}\] \ans Net depth $88.5$ mm; gross $110.7$ mm \ans \[Q=10Ad/(E_aTh)=10\times10\times88.5/(0.8\times14\times20)=39.5\] m$^3$\,h$^{-1}$ ($11.0$ L\,s$^{-1}$) \finans

}{}
\IfFileExists{extra/teny07.tex}{\subsection{Worked solutions on the ten soils (third scenario)}
The same questions are repeated for different working conditions, to show how the answer moves with them.

\exbox{Ten-soil example 7c: Coarse loam (S01)}{Coarse loam: size the irrigation interval and system capacity of a 40-ha maize farm at peak $ET_c=6.5$ mm\,d$^{-1}$, root zone 0.6 m, $p=0.5$, $E_a=0.75$, working 20 h per day.}\sol \givenline{$ET_c=6.5$ mm\,d$^{-1}$ ; $p=0.5$ ; $E_a=0.75$}\ans $TAW=90$ mm, $RAW=45$ mm \ans Interval \[T=\lfloor RAW/ET_c\rfloor=\lfloor6.9\rfloor=6\ \text{d}\] \ans Net depth $T\,ET_c=39.0$ mm; gross $=39.0/0.75=52.0$ mm \ans Capacity \[Q=10Ad/(E_aTh)=10\times40\times39.0/(0.75\times6\times20)=173\] m$^3$\,h$^{-1}$ ($48$ L\,s$^{-1}$) \finans

\exbox{Ten-soil example 7c: Coarse sand (S02)}{Coarse sand: size the irrigation interval and system capacity of a 40-ha maize farm at peak $ET_c=6.5$ mm\,d$^{-1}$, root zone 0.6 m, $p=0.5$, $E_a=0.75$, working 20 h per day.}\sol \givenline{$ET_c=6.5$ mm\,d$^{-1}$ ; $p=0.5$ ; $E_a=0.75$}\ans $TAW=30$ mm, $RAW=15$ mm \ans Interval \[T=\lfloor RAW/ET_c\rfloor=\lfloor2.3\rfloor=2\ \text{d}\] \ans Net depth $T\,ET_c=13.0$ mm; gross $=13.0/0.75=17.3$ mm \ans Capacity \[Q=10Ad/(E_aTh)=10\times40\times13.0/(0.75\times2\times20)=173\] m$^3$\,h$^{-1}$ ($48$ L\,s$^{-1}$) \finans

\exbox{Ten-soil example 7c: Sand (S03)}{Sand: size the irrigation interval and system capacity of a 40-ha maize farm at peak $ET_c=6.5$ mm\,d$^{-1}$, root zone 0.6 m, $p=0.5$, $E_a=0.75$, working 20 h per day.}\sol \givenline{$ET_c=6.5$ mm\,d$^{-1}$ ; $p=0.5$ ; $E_a=0.75$}\ans $TAW=36$ mm, $RAW=18$ mm \ans Interval \[T=\lfloor RAW/ET_c\rfloor=\lfloor2.8\rfloor=2\ \text{d}\] \ans Net depth $T\,ET_c=13.0$ mm; gross $=13.0/0.75=17.3$ mm \ans Capacity \[Q=10Ad/(E_aTh)=10\times40\times13.0/(0.75\times2\times20)=173\] m$^3$\,h$^{-1}$ ($48$ L\,s$^{-1}$) \finans

\exbox{Ten-soil example 7c: Loamy sand (S04)}{Loamy sand: size the irrigation interval and system capacity of a 40-ha maize farm at peak $ET_c=6.5$ mm\,d$^{-1}$, root zone 0.6 m, $p=0.5$, $E_a=0.75$, working 20 h per day.}\sol \givenline{$ET_c=6.5$ mm\,d$^{-1}$ ; $p=0.5$ ; $E_a=0.75$}\ans $TAW=54$ mm, $RAW=27$ mm \ans Interval \[T=\lfloor RAW/ET_c\rfloor=\lfloor4.2\rfloor=4\ \text{d}\] \ans Net depth $T\,ET_c=26.0$ mm; gross $=26.0/0.75=34.7$ mm \ans Capacity \[Q=10Ad/(E_aTh)=10\times40\times26.0/(0.75\times4\times20)=173\] m$^3$\,h$^{-1}$ ($48$ L\,s$^{-1}$) \finans

\exbox{Ten-soil example 7c: Sandy loam (S05)}{Sandy loam: size the irrigation interval and system capacity of a 40-ha maize farm at peak $ET_c=6.5$ mm\,d$^{-1}$, root zone 0.6 m, $p=0.5$, $E_a=0.75$, working 20 h per day.}\sol \givenline{$ET_c=6.5$ mm\,d$^{-1}$ ; $p=0.5$ ; $E_a=0.75$}\ans $TAW=72$ mm, $RAW=36$ mm \ans Interval \[T=\lfloor RAW/ET_c\rfloor=\lfloor5.5\rfloor=5\ \text{d}\] \ans Net depth $T\,ET_c=32.5$ mm; gross $=32.5/0.75=43.3$ mm \ans Capacity \[Q=10Ad/(E_aTh)=10\times40\times32.5/(0.75\times5\times20)=173\] m$^3$\,h$^{-1}$ ($48$ L\,s$^{-1}$) \finans

\exbox{Ten-soil example 7c: Loam (S06)}{Loam: size the irrigation interval and system capacity of a 40-ha maize farm at peak $ET_c=6.5$ mm\,d$^{-1}$, root zone 0.6 m, $p=0.5$, $E_a=0.75$, working 20 h per day.}\sol \givenline{$ET_c=6.5$ mm\,d$^{-1}$ ; $p=0.5$ ; $E_a=0.75$}\ans $TAW=102$ mm, $RAW=51$ mm \ans Interval \[T=\lfloor RAW/ET_c\rfloor=\lfloor7.8\rfloor=7\ \text{d}\] \ans Net depth $T\,ET_c=45.5$ mm; gross $=45.5/0.75=60.7$ mm \ans Capacity \[Q=10Ad/(E_aTh)=10\times40\times45.5/(0.75\times7\times20)=173\] m$^3$\,h$^{-1}$ ($48$ L\,s$^{-1}$) \finans

\exbox{Ten-soil example 7c: Silt loam (S07)}{Silt loam: size the irrigation interval and system capacity of a 40-ha maize farm at peak $ET_c=6.5$ mm\,d$^{-1}$, root zone 0.6 m, $p=0.5$, $E_a=0.75$, working 20 h per day.}\sol \givenline{$ET_c=6.5$ mm\,d$^{-1}$ ; $p=0.5$ ; $E_a=0.75$}\ans $TAW=120$ mm, $RAW=60$ mm \ans Interval \[T=\lfloor RAW/ET_c\rfloor=\lfloor9.2\rfloor=9\ \text{d}\] \ans Net depth $T\,ET_c=58.5$ mm; gross $=58.5/0.75=78.0$ mm \ans Capacity \[Q=10Ad/(E_aTh)=10\times40\times58.5/(0.75\times9\times20)=173\] m$^3$\,h$^{-1}$ ($48$ L\,s$^{-1}$) \finans

\exbox{Ten-soil example 7c: Sandy clay loam (S08)}{Sandy clay loam: size the irrigation interval and system capacity of a 40-ha maize farm at peak $ET_c=6.5$ mm\,d$^{-1}$, root zone 0.6 m, $p=0.5$, $E_a=0.75$, working 20 h per day.}\sol \givenline{$ET_c=6.5$ mm\,d$^{-1}$ ; $p=0.5$ ; $E_a=0.75$}\ans $TAW=84$ mm, $RAW=42$ mm \ans Interval \[T=\lfloor RAW/ET_c\rfloor=\lfloor6.5\rfloor=6\ \text{d}\] \ans Net depth $T\,ET_c=39.0$ mm; gross $=39.0/0.75=52.0$ mm \ans Capacity \[Q=10Ad/(E_aTh)=10\times40\times39.0/(0.75\times6\times20)=173\] m$^3$\,h$^{-1}$ ($48$ L\,s$^{-1}$) \finans

\exbox{Ten-soil example 7c: Clay loam (S09)}{Clay loam: size the irrigation interval and system capacity of a 40-ha maize farm at peak $ET_c=6.5$ mm\,d$^{-1}$, root zone 0.6 m, $p=0.5$, $E_a=0.75$, working 20 h per day.}\sol \givenline{$ET_c=6.5$ mm\,d$^{-1}$ ; $p=0.5$ ; $E_a=0.75$}\ans $TAW=96$ mm, $RAW=48$ mm \ans Interval \[T=\lfloor RAW/ET_c\rfloor=\lfloor7.4\rfloor=7\ \text{d}\] \ans Net depth $T\,ET_c=45.5$ mm; gross $=45.5/0.75=60.7$ mm \ans Capacity \[Q=10Ad/(E_aTh)=10\times40\times45.5/(0.75\times7\times20)=173\] m$^3$\,h$^{-1}$ ($48$ L\,s$^{-1}$) \finans

\exbox{Ten-soil example 7c: Silty clay loam (S10)}{Silty clay loam: size the irrigation interval and system capacity of a 40-ha maize farm at peak $ET_c=6.5$ mm\,d$^{-1}$, root zone 0.6 m, $p=0.5$, $E_a=0.75$, working 20 h per day.}\sol \givenline{$ET_c=6.5$ mm\,d$^{-1}$ ; $p=0.5$ ; $E_a=0.75$}\ans $TAW=102$ mm, $RAW=51$ mm \ans Interval \[T=\lfloor RAW/ET_c\rfloor=\lfloor7.8\rfloor=7\ \text{d}\] \ans Net depth $T\,ET_c=45.5$ mm; gross $=45.5/0.75=60.7$ mm \ans Capacity \[Q=10Ad/(E_aTh)=10\times40\times45.5/(0.75\times7\times20)=173\] m$^3$\,h$^{-1}$ ($48$ L\,s$^{-1}$) \finans

}{}
\IfFileExists{extra/tenby07.tex}{\subsection{Worked solutions on ten crops (hot climate)}
The ten crops of the library differ in their demand, their roots and their sensitivity. The same type of question is answered for each of them at $ET_0=7.5$ mm\,d$^{-1}$; the data are typical values of the crops (Doorenbos--Kassam and FAO 56).

\exbox{Crop example 7c: Maize}{Maize on 10 ha of the loam S06: design the interval, the gross depth and the flow to irrigate in 20 h\,d$^{-1}$ at $E_a=0.80$.}\sol \givenline{$E_a=0.80$}\ans $RAW=94$ mm; \par\indent We know that
\[T=\lfloor94/9.00\rfloor\ti{i}\]
Substituting the values in equation \rfi{i},
\[T=10\ \text{d}\] \ans Net depth $90.0$ mm; gross $112.5$ mm \ans \par\indent We know that
\[Q=10Ad/(E_aTh)\ti{ii}\]
Substituting the values in equation \rfi{ii},
\[Q=10\times10\times90.0/(0.8\times10\times20)=56.2\] m$^3$\,h$^{-1}$ ($15.6$ L\,s$^{-1}$) \finans

\exbox{Crop example 7c: Wheat}{Wheat on 10 ha of the loam S06: design the interval, the gross depth and the flow to irrigate in 20 h\,d$^{-1}$ at $E_a=0.80$.}\sol \givenline{$E_a=0.80$}\ans $RAW=112$ mm; \par\indent We know that
\[T=\lfloor112/8.62\rfloor\ti{i}\]
Substituting the values in equation \rfi{i},
\[T=13\ \text{d}\] \ans Net depth $112.1$ mm; gross $140.2$ mm \ans \par\indent We know that
\[Q=10Ad/(E_aTh)\ti{ii}\]
Substituting the values in equation \rfi{ii},
\[Q=10\times10\times112.1/(0.8\times13\times20)=53.9\] m$^3$\,h$^{-1}$ ($15.0$ L\,s$^{-1}$) \finans

\exbox{Crop example 7c: Cotton}{Cotton on 10 ha of the loam S06: design the interval, the gross depth and the flow to irrigate in 20 h\,d$^{-1}$ at $E_a=0.80$.}\sol \givenline{$E_a=0.80$}\ans $RAW=122$ mm; \par\indent We know that
\[T=\lfloor122/8.62\rfloor\ti{i}\]
Substituting the values in equation \rfi{i},
\[T=14\ \text{d}\] \ans Net depth $120.7$ mm; gross $150.9$ mm \ans \par\indent We know that
\[Q=10Ad/(E_aTh)\ti{ii}\]
Substituting the values in equation \rfi{ii},
\[Q=10\times10\times120.7/(0.8\times14\times20)=53.9\] m$^3$\,h$^{-1}$ ($15.0$ L\,s$^{-1}$) \finans

\exbox{Crop example 7c: Tomato}{Tomato on 10 ha of the loam S06: design the interval, the gross depth and the flow to irrigate in 20 h\,d$^{-1}$ at $E_a=0.80$.}\sol \givenline{$E_a=0.80$}\ans $RAW=48$ mm; $T=\lfloor48/8.62\rfloor=5$ d \ans Net depth $43.1$ mm; gross $53.9$ mm \ans \par\indent We know that
\[Q=10Ad/(E_aTh)\ti{i}\]
Substituting the values in equation \rfi{i},
\[Q=10\times10\times43.1/(0.8\times5\times20)=53.9\] m$^3$\,h$^{-1}$ ($15.0$ L\,s$^{-1}$) \finans

\exbox{Crop example 7c: Potato}{Potato on 10 ha of the loam S06: design the interval, the gross depth and the flow to irrigate in 20 h\,d$^{-1}$ at $E_a=0.80$.}\sol \givenline{$E_a=0.80$}\ans $RAW=30$ mm; $T=\lfloor30/8.62\rfloor=3$ d \ans Net depth $25.9$ mm; gross $32.3$ mm \ans \par\indent We know that
\[Q=10Ad/(E_aTh)\ti{i}\]
Substituting the values in equation \rfi{i},
\[Q=10\times10\times25.9/(0.8\times3\times20)=53.9\] m$^3$\,h$^{-1}$ ($15.0$ L\,s$^{-1}$) \finans

\exbox{Crop example 7c: Soybean}{Soybean on 10 ha of the loam S06: design the interval, the gross depth and the flow to irrigate in 20 h\,d$^{-1}$ at $E_a=0.80$.}\sol \givenline{$E_a=0.80$}\ans $RAW=51$ mm; $T=\lfloor51/8.62\rfloor=5$ d \ans Net depth $43.1$ mm; gross $53.9$ mm \ans \par\indent We know that
\[Q=10Ad/(E_aTh)\ti{i}\]
Substituting the values in equation \rfi{i},
\[Q=10\times10\times43.1/(0.8\times5\times20)=53.9\] m$^3$\,h$^{-1}$ ($15.0$ L\,s$^{-1}$) \finans

\exbox{Crop example 7c: Sugar beet}{Sugar beet on 10 ha of the loam S06: design the interval, the gross depth and the flow to irrigate in 20 h\,d$^{-1}$ at $E_a=0.80$.}\sol \givenline{$E_a=0.80$}\ans $RAW=84$ mm; $T=\lfloor84/9.00\rfloor=9$ d \ans Net depth $81.0$ mm; gross $101.2$ mm \ans \par\indent We know that
\[Q=10Ad/(E_aTh)\ti{i}\]
Substituting the values in equation \rfi{i},
\[Q=10\times10\times81.0/(0.8\times9\times20)=56.2\] m$^3$\,h$^{-1}$ ($15.6$ L\,s$^{-1}$) \finans

\exbox{Crop example 7c: Sunflower}{Sunflower on 10 ha of the loam S06: design the interval, the gross depth and the flow to irrigate in 20 h\,d$^{-1}$ at $E_a=0.80$.}\sol \givenline{$E_a=0.80$}\ans $RAW=61$ mm; $T=\lfloor61/8.25\rfloor=7$ d \ans Net depth $57.8$ mm; gross $72.2$ mm \ans \par\indent We know that
\[Q=10Ad/(E_aTh)\ti{i}\]
Substituting the values in equation \rfi{i},
\[Q=10\times10\times57.8/(0.8\times7\times20)=51.6\] m$^3$\,h$^{-1}$ ($14.3$ L\,s$^{-1}$) \finans

\exbox{Crop example 7c: Onion}{Onion on 10 ha of the loam S06: design the interval, the gross depth and the flow to irrigate in 20 h\,d$^{-1}$ at $E_a=0.80$.}\sol \givenline{$E_a=0.80$}\ans $RAW=18$ mm; $T=\lfloor18/7.88\rfloor=2$ d \ans Net depth $15.8$ mm; gross $19.7$ mm \ans \par\indent We know that
\[Q=10Ad/(E_aTh)\ti{i}\]
Substituting the values in equation \rfi{i},
\[Q=10\times10\times15.8/(0.8\times2\times20)=49.2\] m$^3$\,h$^{-1}$ ($13.7$ L\,s$^{-1}$) \finans

\exbox{Crop example 7c: Alfalfa}{Alfalfa on 10 ha of the loam S06: design the interval, the gross depth and the flow to irrigate in 20 h\,d$^{-1}$ at $E_a=0.80$.}\sol \givenline{$E_a=0.80$}\ans $RAW=94$ mm; \par\indent We know that
\[T=\lfloor94/8.62\rfloor\ti{i}\]
Substituting the values in equation \rfi{i},
\[T=10\ \text{d}\] \ans Net depth $86.2$ mm; gross $107.8$ mm \ans \par\indent We know that
\[Q=10Ad/(E_aTh)\ti{ii}\]
Substituting the values in equation \rfi{ii},
\[Q=10\times10\times86.2/(0.8\times10\times20)=53.9\] m$^3$\,h$^{-1}$ ($15.0$ L\,s$^{-1}$) \finans

}{}
\IfFileExists{extra/tenx07.tex}{\subsection{Worked solutions on the ten soils (second scenario)}
The same questions are repeated for different working conditions, to show how the answer moves with them.

\exbox{Ten-soil example 7b: Coarse loam (S01)}{Coarse loam: size the irrigation interval and system capacity of a 40-ha maize farm at peak $ET_c=6.5$ mm\,d$^{-1}$, root zone 0.6 m, $p=0.5$, $E_a=0.75$, working 20 h per day.}\sol \givenline{$ET_c=6.5$ mm\,d$^{-1}$ ; $p=0.5$ ; $E_a=0.75$}\ans $TAW=90$ mm, $RAW=45$ mm \ans Interval \[T=\lfloor RAW/ET_c\rfloor=\lfloor6.9\rfloor=6\ \text{d}\] \ans Net depth $T\,ET_c=39.0$ mm; gross $=39.0/0.75=52.0$ mm \ans Capacity \[Q=10Ad/(E_aTh)=10\times40\times39.0/(0.75\times6\times20)=173\] m$^3$\,h$^{-1}$ ($48$ L\,s$^{-1}$) \finans

\exbox{Ten-soil example 7b: Coarse sand (S02)}{Coarse sand: size the irrigation interval and system capacity of a 40-ha maize farm at peak $ET_c=6.5$ mm\,d$^{-1}$, root zone 0.6 m, $p=0.5$, $E_a=0.75$, working 20 h per day.}\sol \givenline{$ET_c=6.5$ mm\,d$^{-1}$ ; $p=0.5$ ; $E_a=0.75$}\ans $TAW=30$ mm, $RAW=15$ mm \ans Interval \[T=\lfloor RAW/ET_c\rfloor=\lfloor2.3\rfloor=2\ \text{d}\] \ans Net depth $T\,ET_c=13.0$ mm; gross $=13.0/0.75=17.3$ mm \ans Capacity \[Q=10Ad/(E_aTh)=10\times40\times13.0/(0.75\times2\times20)=173\] m$^3$\,h$^{-1}$ ($48$ L\,s$^{-1}$) \finans

\exbox{Ten-soil example 7b: Sand (S03)}{Sand: size the irrigation interval and system capacity of a 40-ha maize farm at peak $ET_c=6.5$ mm\,d$^{-1}$, root zone 0.6 m, $p=0.5$, $E_a=0.75$, working 20 h per day.}\sol \givenline{$ET_c=6.5$ mm\,d$^{-1}$ ; $p=0.5$ ; $E_a=0.75$}\ans $TAW=36$ mm, $RAW=18$ mm \ans Interval \[T=\lfloor RAW/ET_c\rfloor=\lfloor2.8\rfloor=2\ \text{d}\] \ans Net depth $T\,ET_c=13.0$ mm; gross $=13.0/0.75=17.3$ mm \ans Capacity \[Q=10Ad/(E_aTh)=10\times40\times13.0/(0.75\times2\times20)=173\] m$^3$\,h$^{-1}$ ($48$ L\,s$^{-1}$) \finans

\exbox{Ten-soil example 7b: Loamy sand (S04)}{Loamy sand: size the irrigation interval and system capacity of a 40-ha maize farm at peak $ET_c=6.5$ mm\,d$^{-1}$, root zone 0.6 m, $p=0.5$, $E_a=0.75$, working 20 h per day.}\sol \givenline{$ET_c=6.5$ mm\,d$^{-1}$ ; $p=0.5$ ; $E_a=0.75$}\ans $TAW=54$ mm, $RAW=27$ mm \ans Interval \[T=\lfloor RAW/ET_c\rfloor=\lfloor4.2\rfloor=4\ \text{d}\] \ans Net depth $T\,ET_c=26.0$ mm; gross $=26.0/0.75=34.7$ mm \ans Capacity \[Q=10Ad/(E_aTh)=10\times40\times26.0/(0.75\times4\times20)=173\] m$^3$\,h$^{-1}$ ($48$ L\,s$^{-1}$) \finans

\exbox{Ten-soil example 7b: Sandy loam (S05)}{Sandy loam: size the irrigation interval and system capacity of a 40-ha maize farm at peak $ET_c=6.5$ mm\,d$^{-1}$, root zone 0.6 m, $p=0.5$, $E_a=0.75$, working 20 h per day.}\sol \givenline{$ET_c=6.5$ mm\,d$^{-1}$ ; $p=0.5$ ; $E_a=0.75$}\ans $TAW=72$ mm, $RAW=36$ mm \ans Interval \[T=\lfloor RAW/ET_c\rfloor=\lfloor5.5\rfloor=5\ \text{d}\] \ans Net depth $T\,ET_c=32.5$ mm; gross $=32.5/0.75=43.3$ mm \ans Capacity \[Q=10Ad/(E_aTh)=10\times40\times32.5/(0.75\times5\times20)=173\] m$^3$\,h$^{-1}$ ($48$ L\,s$^{-1}$) \finans

\exbox{Ten-soil example 7b: Loam (S06)}{Loam: size the irrigation interval and system capacity of a 40-ha maize farm at peak $ET_c=6.5$ mm\,d$^{-1}$, root zone 0.6 m, $p=0.5$, $E_a=0.75$, working 20 h per day.}\sol \givenline{$ET_c=6.5$ mm\,d$^{-1}$ ; $p=0.5$ ; $E_a=0.75$}\ans $TAW=102$ mm, $RAW=51$ mm \ans Interval \[T=\lfloor RAW/ET_c\rfloor=\lfloor7.8\rfloor=7\ \text{d}\] \ans Net depth $T\,ET_c=45.5$ mm; gross $=45.5/0.75=60.7$ mm \ans Capacity \[Q=10Ad/(E_aTh)=10\times40\times45.5/(0.75\times7\times20)=173\] m$^3$\,h$^{-1}$ ($48$ L\,s$^{-1}$) \finans

\exbox{Ten-soil example 7b: Silt loam (S07)}{Silt loam: size the irrigation interval and system capacity of a 40-ha maize farm at peak $ET_c=6.5$ mm\,d$^{-1}$, root zone 0.6 m, $p=0.5$, $E_a=0.75$, working 20 h per day.}\sol \givenline{$ET_c=6.5$ mm\,d$^{-1}$ ; $p=0.5$ ; $E_a=0.75$}\ans $TAW=120$ mm, $RAW=60$ mm \ans Interval \[T=\lfloor RAW/ET_c\rfloor=\lfloor9.2\rfloor=9\ \text{d}\] \ans Net depth $T\,ET_c=58.5$ mm; gross $=58.5/0.75=78.0$ mm \ans Capacity \[Q=10Ad/(E_aTh)=10\times40\times58.5/(0.75\times9\times20)=173\] m$^3$\,h$^{-1}$ ($48$ L\,s$^{-1}$) \finans

\exbox{Ten-soil example 7b: Sandy clay loam (S08)}{Sandy clay loam: size the irrigation interval and system capacity of a 40-ha maize farm at peak $ET_c=6.5$ mm\,d$^{-1}$, root zone 0.6 m, $p=0.5$, $E_a=0.75$, working 20 h per day.}\sol \givenline{$ET_c=6.5$ mm\,d$^{-1}$ ; $p=0.5$ ; $E_a=0.75$}\ans $TAW=84$ mm, $RAW=42$ mm \ans Interval \[T=\lfloor RAW/ET_c\rfloor=\lfloor6.5\rfloor=6\ \text{d}\] \ans Net depth $T\,ET_c=39.0$ mm; gross $=39.0/0.75=52.0$ mm \ans Capacity \[Q=10Ad/(E_aTh)=10\times40\times39.0/(0.75\times6\times20)=173\] m$^3$\,h$^{-1}$ ($48$ L\,s$^{-1}$) \finans

\exbox{Ten-soil example 7b: Clay loam (S09)}{Clay loam: size the irrigation interval and system capacity of a 40-ha maize farm at peak $ET_c=6.5$ mm\,d$^{-1}$, root zone 0.6 m, $p=0.5$, $E_a=0.75$, working 20 h per day.}\sol \givenline{$ET_c=6.5$ mm\,d$^{-1}$ ; $p=0.5$ ; $E_a=0.75$}\ans $TAW=96$ mm, $RAW=48$ mm \ans Interval \[T=\lfloor RAW/ET_c\rfloor=\lfloor7.4\rfloor=7\ \text{d}\] \ans Net depth $T\,ET_c=45.5$ mm; gross $=45.5/0.75=60.7$ mm \ans Capacity \[Q=10Ad/(E_aTh)=10\times40\times45.5/(0.75\times7\times20)=173\] m$^3$\,h$^{-1}$ ($48$ L\,s$^{-1}$) \finans

\exbox{Ten-soil example 7b: Silty clay loam (S10)}{Silty clay loam: size the irrigation interval and system capacity of a 40-ha maize farm at peak $ET_c=6.5$ mm\,d$^{-1}$, root zone 0.6 m, $p=0.5$, $E_a=0.75$, working 20 h per day.}\sol \givenline{$ET_c=6.5$ mm\,d$^{-1}$ ; $p=0.5$ ; $E_a=0.75$}\ans $TAW=102$ mm, $RAW=51$ mm \ans Interval \[T=\lfloor RAW/ET_c\rfloor=\lfloor7.8\rfloor=7\ \text{d}\] \ans Net depth $T\,ET_c=45.5$ mm; gross $=45.5/0.75=60.7$ mm \ans Capacity \[Q=10Ad/(E_aTh)=10\times40\times45.5/(0.75\times7\times20)=173\] m$^3$\,h$^{-1}$ ($48$ L\,s$^{-1}$) \finans

}{}
\IfFileExists{extra/tenbx07.tex}{\section{Worked solutions on ten crops (cool climate)}
The ten crops of the library differ in their demand, their roots and their sensitivity. The same type of question is answered for each of them at $ET_0=4.0$ mm\,d$^{-1}$; the data are typical values of the crops (Doorenbos--Kassam and FAO 56).

\exbox{Crop example 7b: Maize}{Maize on 10 ha of the loam S06: design the interval, the gross depth and the flow to irrigate in 20 h\,d$^{-1}$ at $E_a=0.80$.}\sol $RAW=94$ mm; $T=\lfloor94/4.80\rfloor=19$ d \ans Net depth $91.2$ mm; gross $114.0$ mm \ans $Q=10Ad/(E_aTh)=10\times10\times91.2/(0.8\times19\times20)=30.0$ m$^3$\,h$^{-1}$ ($8.3$ L\,s$^{-1}$) \ans

\exbox{Crop example 7b: Wheat}{Wheat on 10 ha of the loam S06: design the interval, the gross depth and the flow to irrigate in 20 h\,d$^{-1}$ at $E_a=0.80$.}\sol $RAW=112$ mm; $T=\lfloor112/4.60\rfloor=24$ d \ans Net depth $110.4$ mm; gross $138.0$ mm \ans $Q=10Ad/(E_aTh)=10\times10\times110.4/(0.8\times24\times20)=28.7$ m$^3$\,h$^{-1}$ ($8.0$ L\,s$^{-1}$) \ans

\exbox{Crop example 7b: Cotton}{Cotton on 10 ha of the loam S06: design the interval, the gross depth and the flow to irrigate in 20 h\,d$^{-1}$ at $E_a=0.80$.}\sol $RAW=122$ mm; $T=\lfloor122/4.60\rfloor=26$ d \ans Net depth $119.6$ mm; gross $149.5$ mm \ans $Q=10Ad/(E_aTh)=10\times10\times119.6/(0.8\times26\times20)=28.8$ m$^3$\,h$^{-1}$ ($8.0$ L\,s$^{-1}$) \ans

\exbox{Crop example 7b: Tomato}{Tomato on 10 ha of the loam S06: design the interval, the gross depth and the flow to irrigate in 20 h\,d$^{-1}$ at $E_a=0.80$.}\sol $RAW=48$ mm; $T=\lfloor48/4.60\rfloor=10$ d \ans Net depth $46.0$ mm; gross $57.5$ mm \ans $Q=10Ad/(E_aTh)=10\times10\times46.0/(0.8\times10\times20)=28.8$ m$^3$\,h$^{-1}$ ($8.0$ L\,s$^{-1}$) \ans

\exbox{Crop example 7b: Potato}{Potato on 10 ha of the loam S06: design the interval, the gross depth and the flow to irrigate in 20 h\,d$^{-1}$ at $E_a=0.80$.}\sol $RAW=30$ mm; $T=\lfloor30/4.60\rfloor=6$ d \ans Net depth $27.6$ mm; gross $34.5$ mm \ans $Q=10Ad/(E_aTh)=10\times10\times27.6/(0.8\times6\times20)=28.7$ m$^3$\,h$^{-1}$ ($8.0$ L\,s$^{-1}$) \ans

\exbox{Crop example 7b: Soybean}{Soybean on 10 ha of the loam S06: design the interval, the gross depth and the flow to irrigate in 20 h\,d$^{-1}$ at $E_a=0.80$.}\sol $RAW=51$ mm; $T=\lfloor51/4.60\rfloor=11$ d \ans Net depth $50.6$ mm; gross $63.2$ mm \ans $Q=10Ad/(E_aTh)=10\times10\times50.6/(0.8\times11\times20)=28.7$ m$^3$\,h$^{-1}$ ($8.0$ L\,s$^{-1}$) \ans

\exbox{Crop example 7b: Sugar beet}{Sugar beet on 10 ha of the loam S06: design the interval, the gross depth and the flow to irrigate in 20 h\,d$^{-1}$ at $E_a=0.80$.}\sol $RAW=84$ mm; $T=\lfloor84/4.80\rfloor=17$ d \ans Net depth $81.6$ mm; gross $102.0$ mm \ans $Q=10Ad/(E_aTh)=10\times10\times81.6/(0.8\times17\times20)=30.0$ m$^3$\,h$^{-1}$ ($8.3$ L\,s$^{-1}$) \ans

\exbox{Crop example 7b: Sunflower}{Sunflower on 10 ha of the loam S06: design the interval, the gross depth and the flow to irrigate in 20 h\,d$^{-1}$ at $E_a=0.80$.}\sol $RAW=61$ mm; $T=\lfloor61/4.40\rfloor=13$ d \ans Net depth $57.2$ mm; gross $71.5$ mm \ans $Q=10Ad/(E_aTh)=10\times10\times57.2/(0.8\times13\times20)=27.5$ m$^3$\,h$^{-1}$ ($7.6$ L\,s$^{-1}$) \ans

\exbox{Crop example 7b: Onion}{Onion on 10 ha of the loam S06: design the interval, the gross depth and the flow to irrigate in 20 h\,d$^{-1}$ at $E_a=0.80$.}\sol $RAW=18$ mm; $T=\lfloor18/4.20\rfloor=4$ d \ans Net depth $16.8$ mm; gross $21.0$ mm \ans $Q=10Ad/(E_aTh)=10\times10\times16.8/(0.8\times4\times20)=26.2$ m$^3$\,h$^{-1}$ ($7.3$ L\,s$^{-1}$) \ans

\exbox{Crop example 7b: Alfalfa}{Alfalfa on 10 ha of the loam S06: design the interval, the gross depth and the flow to irrigate in 20 h\,d$^{-1}$ at $E_a=0.80$.}\sol $RAW=94$ mm; $T=\lfloor94/4.60\rfloor=20$ d \ans Net depth $92.0$ mm; gross $115.0$ mm \ans $Q=10Ad/(E_aTh)=10\times10\times92.0/(0.8\times20\times20)=28.8$ m$^3$\,h$^{-1}$ ($8.0$ L\,s$^{-1}$) \ans

}{}
\IfFileExists{extra/ex07.tex}{\backhead{EXERCISES}{khred}
\begin{exlist}
\item Find the capacity of a system for 25 ha, net depth 40 mm, $E_a=0.8$, 8 days, 18 h per day.\exans{87 m$^3$ h$^{-1}$}
\item RAW = 45 mm and $ET_c=6.4$ mm\,d$^{-1}$. Find the interval.\exans{7.0 d}
\item Evaporation from bare soil is 150 mm. A mulch of strength 0.5 and capacity 1.5 mm meets 30 events. Find $E_m$.\exans{120 mm}
\item $\rho_b$ relaxes from 1.20 to 1.50 with $k=0.015$ d$^{-1}$. Find $\rho_b$ at 100 d.\exans{1.433 Mg m$^{-3}$}
\item An ellipsoidal plant has $a=0.4$, $b=0.6$ m. Find its volume.\exans{0.402 m$^3$}
\item With $D_r=70$ mm, TAW = 120, RAW = 60, find $K_s$.\exans{0.83}
\item $P=450$, $ET_0=1\,050$ mm. Find AI and its class.\exans{0.43; semi-arid}
\item Wheat yields 5.2 Mg\,ha$^{-1}$ with $ET_a=380$ mm. Find WP.\exans{1.37 kg m$^{-3}$}
\item Drainage rises 45 mm at 15 mg\,L$^{-1}$ nitrate-N. Find the extra N loss (kg\,ha$^{-1}$).\exans{6.75 kg/ha}
\item A factorial of four factors with 3, 2, 4 and 2 levels: how many runs?\exans{48}
\end{exlist}
}{}
\IfFileExists{extra/dyk07.tex}{\dykband{DO YOU KNOW ?}
\begin{exlist}
\item What is a layer of the application, and which layers compute rather than read?
\item State the system capacity formula and its terms.
\item How does a mulch change evaporation, and why is the saving smaller than the strength suggests?
\item Why does the benefit of tillage on drainage decline with time?
\item Explain the daily bucket in words.
\item How is the aridity index used?
\item What is water productivity?
\item What is the sign of a balancing loop?
\end{exlist}
}{}
\IfFileExists{extra/obj07.tex}{\dykband{OBJECTIVE TYPE QUESTIONS}
\begin{exlist}
\item The system capacity $Q=10Ad/(E_aTh)$ is in\opt{L/s}{m$^3$/h}{mm}{ha}
\item A mulch's strength is\opt{its capacity}{the fraction by which it lowers vapour flux}{its thickness}{its price}
\item The bulk density after tillage\opt{stays low}{relaxes to a higher value}{falls}{is zero}
\item $AI<0.2$ is\opt{humid}{arid}{semi-arid}{sub-humid}
\item WP is yield divided by\opt{rainfall}{$ET_a$}{area}{runoff}
\item A balancing loop has a\opt{positive product of signs}{negative product of signs}{zero signs}{no arrows}
\item Layers that read archived results include\opt{control and mulch}{design only}{none}{phenotyping only}
\item The number of runs of a factorial is the\opt{sum of levels}{product of levels}{maximum level}{mean level}
\item In the daily bucket, deep percolation occurs when\opt{the soil is below wilting}{the root zone is refilled beyond field capacity}{it is dark}{there is no rain}
\item An ellipsoid of semi-axes $a,a,b$ has volume\opt{$\frac43\pi a^2b$}{$\pi a^2b$}{$2\pi ab$}{$\frac13\pi ab$}
\end{exlist}
\objkey{1.~(b)\quad 2.~(b)\quad 3.~(b)\quad 4.~(b)\quad 5.~(b)\quad 6.~(b)\quad 7.~(a)\quad 8.~(b)\quad 9.~(b)\quad 10.~(a)}
}{}
\swaflvideo{07}
